% =====================================================================
%  main.tex — replace ONLY the \begin{abstract}...\end{abstract} block.
%  Everything else in main.tex stays exactly as it is.
%
%  Changed: the final three sentences. The three \ph{X} placeholders are
%  filled, and two sentences are added so the abstract states the trade-off
%  and its significance rather than only the favourable number.
% =====================================================================

\begin{abstract}
Code clone \emph{detection}---deciding whether two programs are semantically equivalent despite syntactic differences---has matured into a rich body of work spanning source- and binary-level analysis. The inverse problem, however, remains largely unexplored: given a program, can we \emph{generate} a functional clone that preserves its observable behavior while presenting a materially different implementation? This capability is valuable for software diversification, dependency modernization (e.g., replacing a vulnerable or slow library with a safer or faster one), and the resilience of critical-infrastructure software. In this paper, we present \sys, an agentic system that, given a source program and a set of requirements, synthesizes a working code clone through the coordinated action of four Large-Language-Model-driven agents: a \emph{library-substitution} agent that swaps third-party dependencies for functionally comparable but more secure or performant alternatives, a \emph{glue-synthesis} agent that adapts the surrounding program to the new dependencies, a \emph{logic-mutation} agent that rewrites first-party application logic into behaviorally equivalent yet textually divergent code, and a \emph{verification} agent that exercises available test cases to confirm equivalence and localize regressions. \sys couples two tiers of verification: a \emph{local} execute-and-repair loop that guards each transformation, and a \emph{global} equivalence oracle that confirms end-to-end equivalence or pinpoints where the clone diverges. We evaluate \sys on twelve programs that estimate factory measurements and control industrial infrastructure, against an unverified single-shot baseline and an ablation without the verification tier. Conventional measures---compilation, test-suite pass rate, oracle confirmation---fail to separate the three configurations, because every subject injects its dependencies and so never executes the migrated code path under test. A blind review of all 108 generated clones by an independent model does separate them: 33\% of the clones \sys accepts carry serious behavioral defects, against 54\% for the single-shot baseline. \sys accepts far fewer clones, however, and therefore delivers slightly fewer sound ones overall; on this corpus the difference is not statistically significant. The ablation performs indistinguishably from the single-shot baseline, indicating that the verification tier rather than the agent decomposition carries the effect. These early results suggest that agentic, verification-driven mutation is a promising route to automated, trustworthy software diversification, and that evaluating it requires oracles that execute the substituted path directly.
\end{abstract}
